\scanpage{235}
\mspage{181}
% archive sheet number: 234
qui refusent de se laisser ``\unsure{toucher|boucher}'' par $\Phi_0$ ou $\Psi_0$, et
pour lesquelles \unsure{même|mêmes} $C$ n'est d'aucun secours,
faute de l'axiome $W(4\mathrm{b})$.

\textcircled{4} \uline{Cas où l'axiome} $W(4\mathrm{b})$ \uline{est satisfait}

Alors les perplexités que je viens de dire disparaissent. On a aussitôt, en vertu de $W(4\mathrm{b})$ :
\[
\left|\!\left|\!\left|\begin{aligned}
&\text{Si }X\in\mathrm{Ob}\,\underline{\mathrm{Fib}}_W\,S,\quad\text{alors }i_X:X\longrightarrow C(X)\quad\text{est dans }W_S^f\\
&\qquad\text{(\uline{en supposant l'axiome $W(4\mathrm{b})$})}
\end{aligned}\right.\right.\right.
\]
et $C$ fournit des quasi-inverses pour les foncteurs
``d'inclusion'' entre localisés fins (par $W_S^f$) envisagés à la
page précédente. \deleted{Il en fournit} un également
pour le foncteur $\lambda$ envisagé plus haut ($(*)$ page 178)
\diagram{\begin{tikzcd}[ampersand replacement=\&,row sep=.15em]
(\underline{\mathrm{Fib}}_W\,S)(W_S^{\deleted{\unsure{f}}})^{-1}\arrow[r]\arrow[d,phantom,"\Vert\ {\mathrm{def}}"] \&
(\mathrm{Cat}_{/S})(W_S^{-1})\arrow[d,phantom,"\deleted{\unsure{\Vert}}\,\unsure{\wr\wr|{??}}"]\\
\mathrm{HOT}_W^{\mathrm{lc}}(S) \& \mathrm{HOT}_W^{\mathrm{lc}!}(S)\;{,}
\end{tikzcd}}
lequel est donc une équivalence. Une autre façon,
plus évidente, de dire la $=$ chose, c'est que sur
\inserted{chacune} des catégories de l'étage les \inserted{deux} localiseurs $W_S^f$ et $W_S$
sont désormais identiques (puisqu'elles sont toutes
contenues dans $\underline{\mathrm{Fib}}_W\,S$, où elles le sont justement par
l'axiome $W(4\mathrm{b})$), donc le foncteur $\lambda$ entre
ces deux localisés (définition originelle de $\lambda$,
dans le cas $\unsure{\mathcal{F}|F}=\underline{\mathrm{Fib}}_W\,S$, cf \unsure{page} 176) est même
le \uline{foncteur identique}, les deux catégories \inserted{localisées} étant \uline{égales} :
\[
\mathrm{HOT}_W^{\mathrm{lc}}(S)=\mathrm{HOT}_W^{\mathrm{lc}!}(S)\qquad !
\]
Ainsi le diagramme $(*)$ p. 178 prend la forme
\diagram{\begin{tikzcd}[ampersand replacement=\&,column sep=small]
 \&\mathrm{HOT}_W^{\mathrm{lc}}(S)\arrow[dl,"i"]\arrow[dr,"j"'] \&\\
 (*)\quad\mathrm{HOT}_W(S)\arrow[ur,bend left=17,"\rho"] \& \&\mathrm{HOT}_W^{\unsure{0|o}}(S)\arrow[ul,bend right=17,"\sigma"'] .
\end{tikzcd}}
